% !Mode:: "TeX:UTF-8"
%  Authors: 张井   Jing Zhang: prayever@gmail.com     天津大学2010级管理与经济学部信息管理与信息系统专业硕士生
%           余蓝涛 Lantao Yu: lantaoyu1991@gmail.com  天津大学2008级精密仪器与光电子工程学院测控技术与仪器专业本科生

%%%%%%%%%% Package %%%%%%%%%%%%
\usepackage{graphicx}                       % 支持插图处理
% \usepackage[a4paper,text={146.4true mm,239.2 true mm},top= 25.4true mm, bottom= 25.4true mm, left=31.7 true mm,head=6true mm,headsep=6.5true mm,foot=16.5true mm]{geometry}
\usepackage[a4paper,top=25.4mm, bottom=25.4mm, left=31.7mm, right=31.7mm, head=6true mm,headsep=6.5true mm,foot=17.5mm]{geometry}
                                            % 支持版面尺寸设置
\usepackage[squaren]{SIunits}               % 支持国际标准单位

\usepackage{titlesec}                       % 控制标题的宏包
\usepackage{titletoc}                       % 控制目录的宏包
\usepackage{fancyhdr}                       % fancyhdr宏包 支持页眉和页脚的相关定义
\usepackage{color}                          % 支持彩色
\usepackage{amsmath}                        % AMSLaTeX宏包 用来排出更加漂亮的公式
\usepackage{amssymb}                        % 数学符号生成命令
\usepackage[below]{placeins}    %允许上一个section的浮动图形出现在下一个section的开始部分,还提供\FloatBarrier命令,使所有未处理的浮动图形立即被处理
\usepackage{multirow}                       % 使用Multirow宏包，使得表格可以合并多个row格
\usepackage{booktabs}                       % 表格，横的粗线；\specialrule{1pt}{0pt}{0pt}
\usepackage{longtable}                      % 支持跨页的表格。
\usepackage{tabularx}                       % 自动设置表格的列宽
% \usepackage{subfigure}                      % 支持子图（迁移后报告未使用，改用 caption 宏包）
\usepackage{float}                          % 支持 [H] 固定位置图表
\usepackage{caption}                        % 图表题注格式（替代 ccaption）
\usepackage{etoolbox}                       % 提供 \AtBeginEnvironment 等
\usepackage{eso-pic}                        % 提供 \AddToShipoutPictureBG（整页背景封面）
\usepackage{lastpage}                       % 提供 \pageref{LastPage}（页脚“共 N 页”）
\usepackage[sort&compress,numbers]{natbib}  % 支持引用缩写的宏包
\usepackage{enumitem}                       % 使用enumitem宏包,改变列表项的格式
\usepackage{calc}                           % 长度可以用+ - * / 进行计算
\usepackage{bm}                             % 处理数学公式中的黑斜体的宏包
\usepackage[amsmath,thmmarks,hyperref]{ntheorem}  % 定理类环境宏包，其中 amsmath 选项用来兼容 AMS LaTeX 的宏包
\usepackage{indentfirst}                    % 首行缩进宏包

% \usepackage{fancybox} 

%\usepackage{hypbmsec}                      % 用来控制书签中标题显示内容
\newcommand{\tabincell}[2]{\begin{tabular}{@{}#1@{}}#2\end{tabular}}
\usepackage[table]{xcolor}                  % 提供 \rowcolor / \rowcolors 等表格底色命令
%支持代码环境
\usepackage{listings}
\lstset{numbers=left,
language=[ANSI]{C},
numberstyle=\tiny,
extendedchars=false,
showstringspaces=false,
breakatwhitespace=false,
breaklines=true,
captionpos=b,
keywordstyle=\color{blue!70},
commentstyle=\color{red!50!green!50!blue!50},
frame=shadowbox,
rulesepcolor=\color{red!20!green!20!blue!20}
}
%支持算法环境
\usepackage[boxed,ruled,lined]{algorithm2e}
\usepackage{algorithmic}

\usepackage{array}
\newcommand{\PreserveBackslash}[1]{\let\temp=\\#1\let\\=\temp}
\newcolumntype{C}[1]{>{\PreserveBackslash\centering}p{#1}}
\newcolumntype{R}[1]{>{\PreserveBackslash\raggedleft}p{#1}}
\newcolumntype{L}[1]{>{\PreserveBackslash\raggedright}p{#1}}

% 生成有书签的 pdf。使用 xelatex（xeCJK）编译，无需 cmap。
\usepackage{hyperref}
\hypersetup{
    unicode,
    pdfborder={0 0 0},
}
% \usepackage[pdftex,unicode,
%             CJKbookmarks=true,
%             bookmarksnumbered=true,
%             bookmarksopen=true,
%             colorlinks=false,
%             pdfborder={0 0 0},
%             citecolor=blue,
%             linkcolor=red,
%             anchorcolor=green,
%             urlcolor=blue,
%             breaklinks=true
%             ]{hyperref}

